\section{Exact flattening at several scales}\label{gate:multi-v2}

We work first on an arbitrary finite set of $N$ targets. A trial proposes a
target $s$ and decoration $w$, or returns null. Let $Q(s,w)$ be its
forward atom, and $R_s(w)$ the atom in a normalized reverse experiment
at $s$. Both include exactly the stated random labels, marks, and
sampling choices of the construction; internal runtimes are not part
of the decoration. We discard forward atoms outside reverse support
and require $R_s(w)>0\Rightarrow Q(s,w)>0$. Write
\[
 T(s,w)=\frac{R_s(w)}{Q(s,w)},\qquad z(s,w)=\frac{H}{T(s,w)} .
\]
The deterministic positive rational $H$ is common to all targets.
The previous estimates supply two quite different inequalities:
\begin{equation}\label{gate:hyp-v2}
 \mathbb E_{R_s}\log z\ge0\quad(\text{every }s),\qquad H/N=O(1).
\end{equation}
The procedure below does not evaluate either expectation or know $N$.
It also works with a larger $H/N$. Fix a known integer $M\ge2$
with $\log z\le M$ on every reverse history; no corresponding
pointwise lower bound on $\log z$ is needed. For our proposal a
polynomial $M$ suffices. For example, conditional on the public word
and order (whose probabilities cancel), the reverse experiment takes
$O(n)$ matching or coloring marks, each with atom at least
$1/[(n!)^{O(1)} 2^{O(n)}]$, and $O(n)$ independent binary packing
choices. A capped cubic coloring atom satisfies this same crude
bound. Eligible forward probabilities are at most one, and
$\log H=O(n^2\log n)$. We can use $M=n^4$ for all sufficiently large
$n$, with no unknown coefficient needed in this choice.

\begin{lemma}[Multiscale exact flattening]\label{gate:multiscale-v2}
Suppose a forward trial followed by its soft coin of bias $T/(T+a)$,
and a fresh reverse trial at a prescribed $s$ followed by its soft
coin of bias $a/(a+T)$, each cost at most $W$ in expectation,
uniformly for $a=H/2^j$, $0\le j\le\lceil\log_2M\rceil$.
Assume the mean-log inequality in \eqref{gate:hyp-v2}. There is an
exactly uniform sampler of expected total cost
\[
 O\!\left((W+\log^2(2M))(1+H/N)\log^2(2M)\right)
\]
apart from preparing the common inputs. In particular a polynomial
bound on $M$ does not cost a power of $n$.
\end{lemma}

\begin{proof}
Put $J=\lceil\log_2M\rceil$, $m_j=2^j$, $m=m_J$, $u=1/64$, and
\[
 f_j(s)=\mathbb E_{R_s}\frac{z}{1+z/m_j},\qquad
 g_j(s)=f_j(s)/m_j,\qquad r_j(s)=\frac{u}{u+f_j(s)} .
\]
On scale $H/m_j$ the soft forward proposal gives exactly $f_j(s)/H$
mass at $s$; this is the atom identity
$Q T/(T+a)=R_s/(T+a)$. One fresh soft reverse experiment is a
$g_j(s)$-coin. Define $q_s=\mathbb E_{R_s}(1/T)$ only for the proof.
Then $\sum_s q_s\le1$, $\mathbb E_{R_s}z=Hq_s$, and
$0<f_j(s)\le\min(m_j,Hq_s)$. Jensen also gives $Hq_s\ge1$.

For $v\ge2$ and $t\le v$ we have
\[
 \frac{e^t}{1+e^t/v}\ge \frac{1+t}{3}.
\]
When $t\le\log v$ use $e^t/2$ and $e^t\ge1+t$ (a negative right
side is automatic); otherwise use $v/2\ge(1+t)/3$. Thus
$f_J(s)\ge1/3$ for every $s$.

Choose $j$ uniformly in $\{0,\ldots,J\}$ and run its soft forward
proposal. Upon receiving an $s$, require fresh independent $r_i(s)$
tests to pass for $i=0,\ldots,j-1$, stopping at the first failure.
For $j<J$ require one more $r_j$ test. For $j=J$ instead require a
test of bias $u/f_J(s)$. Output $s$ iff all required tests passed.

To implement each $r_i$, repeatedly pick a fair side, try its
$u/m_i$ or $g_i$ coin, and return which side first succeeded. This
costs $O(m_i)$ coin calls uniformly. At the final scale set $c=u/m$
and use the same race to supply $p=c/(c+g_J)$; $p\le3u/(1+3u)<1/5$.
Draw $k\ge1$ with probability $2^{-k}$ and accept iff $k$ independent
$2p$-coins pass. The resulting bias is
$\sum_{k\ge1}p^k=p/(1-p)=u/f_J$. The bounded linear factory of
Huber \cite{Huber} with this fixed slack needs bounded expected input
calls; it uses only fair bits and rational flips. To see bit overhead
can also be bounded at fixed slack, use the doubling-stage version:
at stage $j$ the barrier is $32\,2^j$, relative slack $2^{-j-1}$,
and the amplifier stays between $2$ and $2e^{1/2}$. The exit walk has
mean drift bounded away from zero for $p\le1/5$ and expected length
$O(2^j)$ (the geometric overshoot has bounded mean); its survival
factor to the next stage is at most
$(1+2^{-j-2})^{-32\,2^j}\le e^{-6}$. Rational parameters at this
stage have $O((j+1)^2)$ bits. Implement barrier powers by sequential
rational coins; the geometrically summable stage costs bound auxiliary
bit work. Thus a final test, if invoked, costs $O(m(W+\log^2(2M)))$.

Before the uniform level prior, the output mass at level $j<J$ is
\[
 (f_j/H)\prod_{i\le j}r_i
     =(u/H)(1-r_j)\prod_{i<j}r_i.
\]
At level $J$ it is $(u/H)\prod_{i<J}r_i$. Summing telescopes:
each target receives $u/[H(J+1)]$, exactly.

It remains to charge expensive tests, not simply multiply by their
worst-case cost. Let $B=\{s:Hq_s>m^4\}$, a proof-only set of at most
$H/m^4$ targets. At any level the first-gate mass of $B$ is at most
$(m/H)|B|\le m^{-3}$. Even running every test thereafter costs
$O(m(W+\log^2(2M)))$ on these targets. On all other $s$ put
$K=\log(8m^4)$. Since $(t-K)_+\le e^{t-K}$,
\[
 \mathbb E_{R_s}\min(\log z,K)\ge -e^{-K}Hq_s\ge-1/8 .
\]
Monotonicity and the scalar inequality applied at
$\min(\log z,K)$ give $f_i(s)\ge7/24>1/4$ when $m_i\ge\max(2,K)$.
Then $r_i\le1/17$. Up to the first such $i_0$ the sum of the
worst-case test costs is $O((W+\log^2(2M))\log(2m))$.
Thereafter reaching $i$ costs a probability factor at most
$17^{-(i-i_0)}$; its conditional work grows only as $2^i$.
This bound includes the final inversion. If there is no $i_0$ then
$m<K$ and the same cost bound holds. The first-gate masses form a
subprobability, so these conditional estimates prove the same bound
per complete trial. Independent restarts have a.s. finite stopping
time; the mean number is $H(J+1)/(uN)$. Renewal gives the stated
total expectation even when cost correlates with the success of an
individual trial.
\end{proof}

\subsection{The soft coin, with all matching probes charged}

\begin{lemma}[Binary-prefix permanent gate]\label{gate:binary-v2}
For the decorated row, weighted-fill and cubic proposals in this
paper, both soft coins in Lemma~\ref{gate:multiscale-v2} can be
implemented in expected bit work
$\widetilde O(n^4P(n)+n^{12})$. Here $P(n)$ bounds the expected
bit cost of a fresh uniform perfect matching on every feasible
\emph{unweighted} bipartite board of side at most $n$. This generous
ancillary allowance need not be optimized. No Latin-square completion
oracle, exact permanent evaluation on the ordinary path, or
exponential-time path enumeration is assumed.
\end{lemma}

\begin{proof}
After null histories have been discarded, $T$ is a known positive
rational multiplier times products and reciprocals of ordinary
permanents and positive coin biases. For a balanced weighted fill,
divide the entry-weight product by its explicit common upper
bound from the free-box intersections; simultaneous attainment of that
bound is not required. The resulting weight of a
uniform permissible matching is a rational in $[e^{-4},1]$.
Rejecting by this weight samples the weighted matching, and its
normalizer is the unweighted permanent times a known factor times
the mean of this coin. The occupancy-range argument that gives this
uniform lower bound is important: repeated sampling is not assumed
cheap merely because the individual entries are positive. Imbalanced
compositions use the ordinary fill. A capped cubic block has the
additional factor $J(T_3)=A\per(T_3)\theta$, where $A=n$ suffices and
$\theta=\mathbb E_M\min(1,2^{K(M)}/A)\ge1/A$ for uniform $M$ on
the cubic board. The relevant $K(M)$ is found by following degree-two
cycles. There are only a bounded number of these latter coin factors.
Neither the (possibly very large) individual number of cycles nor a
sum over all perfect matchings is computed to obtain $\theta$.

For each ordinary permanent factor draw an independent uniform
matching called its anchor. Traverse the left vertices in a fixed
order. At a vertex split the list of possible column indices by the
bits of their index in a fixed, padded binary encoding. Reveal the
anchor's next branch bit. Its conditional probability $p$ can be
probed by taking a fresh uniform matching on the current minor with
that row restricted to the branches already revealed. After the
column is fixed, delete that row and column. The product of the
encountered $p$'s equals $1/\per B$ for \emph{each} anchor on board
$B$. At any visited site, averaging over the anchor's next choice
before observing it gives $\mathbb E(1/p)\le2$, the number of
nonempty children. This statement is not claimed conditional on a
complete rare anchor. Use fixed binary depths, even when one child is
forced. Over $O(n)$ boards and the scalar coins above, there are
$D_0=O(n^2\log n)$ sites. Frequencies from a single matching are not
reused as independent probes.

Fix $k=D_0+2$, or use a deterministic upper bound on $D_0$.
At every site write $T_j'$ for the number of independent $p$-coin
flips to obtain $j$ successes. For a required factor $1/p$ take
$T_k'/k$; for a required $p$ take $(k-1)/(T_k'-1)$.
Their means are $1/p$ and $p$, respectively. Their size-biased
versions are $(T_{k+1}'-1)/k$ and $(k-1)/T_{k-1}'$.
These four identities follow directly from
$\Pr(T_j'=t)=\binom{t-1}{j-1}p^j(1-p)^{t-j}$, including $p=1$.
Their relative second moments before size bias are at most
$1+1/k$ and $(k-1)/(k-2)$, the latter by replacing
$(T_k'-1)^{-2}$ by $[(T_k'-1)(T_k'-2)]^{-1}$. Taking independent
factors times the known multiplier yields $Y>0$ with
$\mathbb E Y=T$ and $\mathbb E Y^2\le eT^2$. Both statements hold
conditional on all anchors. The size-biased law $Y^*$ therefore
retains the original anchor law: replace each local factor by its
displayed biased version, independently. Purely deterministic
ratios can of course be sampled without anchors.

Let $l=\mathbb E\min(1,Y/a)$ and $h=\mathbb E\min(1,a/Y^*)$.
Then $Th=al$. The second-moment bound gives
$l\ge(8e)^{-1}\min(1,T/a)$ by Paley--Zygmund. Thus $l+h\ge(8e)^{-1}$
and an independent fair-side race between $l$ and $h$ has bounded
expected rounds. The winning side gives, respectively, $T/(T+a)$
or $a/(T+a)$. This exactness uses neither a logarithmic estimate of
$T$ nor exact knowledge of any of its unknown factors.

A full race round needs $O(kD_0)$ ordinary matching probes in
expectation by the binary-branch identity. The bounded number of
cubic scalar sites add $O(kn)$ and the bounded-away-from-zero fill
sites are covered by $O(kD_0)$. To obtain the decorated forward or
reverse history itself uses only $O(n)$ ordinary draws and at most
$O(n)$ expected draws for the capped cubic rejection, plus its
constant-factor weighted-fill rejections. An independent round may
be long and may win, but renewal with an unconditional constant
success bound still gives $\widetilde O(n^4)$ probes in total.

Here are generous sequential bit allowances for auxiliary work.
A chosen polygon word has fixed length; enumerate its $n$ rooted
candidates, not all assignments of its vertex labels. The composition
normalizer at a $b$-block is a convolution table with $O(bn^2)$
integer operations on $\widetilde O(bn)$-bit entries after factorial
clearing; all block lengths sum to at most $n$. The cubic
double-color histogram has four free counts of range $n$ at each of
$n$ levels. Its $O(n^5)$ transitions can use streaming scans and
fixed-offset shifts of tables in lexicographic order; a constant
number of sorted streams suffices for the six possible vertex types.
Table entries and terminal factorial weights have
$\widetilde O(n)$ bits, and sequential backward sampling may rescan a
full table at each level. The degree-two overflow pair uses the
ordinary integer cycle-index table. These arithmetic costs fit
$\widetilde O(n^9)$ even with schoolbook products and comparisons.

Explicit normalized weighted-fill coins use at most $O(b)$ rational
factor powers, with exponents at most $n$, so their numerator and
denominator have $O(bn\log n)$ bits. They can be computed by repeated
squaring and schoolbook multiplication in $\widetilde O(n^5)$ bit
operations per probe, even using $b\le n$ and scanning every filled
cell for each future free box. The full known path prefactor has
$\widetilde O(n^3)$ bits by summing over the $O(n)$ physical slots;
including $H$ and its dyadic rescalings does not change this bound.
Feasibility, copying a board, and validating a probe cost a polynomial
scan smaller than these allowances. All probe probabilities at
anchor sites are at least $1/n!$. Uniformly in their realized
values, the uncapped negative-binomial counters have
$\mathbb E\log^2(1+T_j')=O((n\log n+\log k)^2)$ by their geometric
tails or Markov's inequality at exponentially increasing multiples
of their mean. For increment costs one may use size bias and
$\mathbb E[T_j'\log(1+T_j')\mid p]\le
(j/p)\log(1+(j+1)/p)$. Thus the unreduced product's squared bitlength
has expectation $\widetilde O(n^6)$, and even $D_0$ schoolbook
multiplications at that bound are harmless here. All rational
choices can be driven by fair bits by sampling a uniform integer
below their denominator. These bounds prove the displayed allowance;
they are bounds on expected total work, not RAM arithmetic counts.
\end{proof}

With $M$ polynomial, $H/N=O(1)$ and the conservative matching
adapter, the preceding two lemmas give the theorem's total bit
bound after adding the one-time scale preparation. An explicit
constant times the calibrated scale may be replaced by an extra
factor $\log(2n)$ for all sufficiently large orders if its implicit
absolute constant is not numerically evaluated; this changes only
logarithmic factors in the bound. The large-order inequalities need
not hold on an unspecified small input. Fix a sufficiently large
admission threshold and use exact exhaustive uniform generation
below it, including $n=1$. This finite modification does not affect
the exponent. It is not a claim that a tested, practical numerical
admission cutoff or a scalable implementation is supplied.
